\section{问题三：多预报发布下的购电调整}

\subsection{信息结构与因果预测}
问题三需要区分预报发布时间和电池控制时刻：购电合同每天只在0:00、6:00、12:00、
18:00制定或调整，电池动作每十分钟重算。将自然日分为
$I_k=[\tau_k,\tau_{k+1})$，$k=0,\ldots,143$，功率按右端点对齐，电量为功率的$1/6$。
在时刻$t$，计划器只读$\mathcal I_t=\{x:\operatorname{available\_at}(x)\le t\}$；
当格实际值在结束后才用于回放结算，不能提前传入计划器。

负载预测使用四周同一时隙的加权基线，并用已实现残差拟合无截距AR(1)：
\begin{align}
 m_u&=\frac{8L_{u-7\mathrm d}+4L_{u-14\mathrm d}+2L_{u-21\mathrm d}+L_{u-28\mathrm d}}{15},\\
 r_u&=L_u-m_u,\qquad
 \widehat L_{t+h}=\max\{0,m_{t+h}+\phi_t^h r_t\}.
\end{align}
其中$h$为十分钟步数，$\phi_t$由截至$t$的历史估计并截断至$[0,0.999]$。
每次发布冻结未来30小时版本，中间MPC只从最近版本截取24小时，不重新拟合。

对相同目标时刻$u$的全部可见光伏版本$j\in\mathcal V_t(u)$，以预报提前量$h_j$
分组，用截至$t$已经实现的误差计算MAE，组合为
\begin{equation}
 \widehat P_{t,u}=\sum_jw_{j,t}P_{j,u},\qquad
 w_{j,t}=\frac{(\operatorname{MAE}_{h_j}(t)+1)^{-2}}
 {\sum_{a\in\mathcal V_t(u)}(\operatorname{MAE}_{h_a}(t)+1)^{-2}}.
\end{equation}
MAE和平滑常数1的单位均为kW。唯一小时序列线性插值到十分钟右端点，当前左边界
使用刚结束区间的实际光伏。窗口超出该版覆盖的后缀，按过去1至7天同一时隙的实际值
补齐，权重$a_j=2^{-j/2}/\sum_{b=1}^7 2^{-b/2}$，半衰期2天。
后发布版本不覆盖历史决策，也不把未来实际值当作补齐来源。

\subsection{合同版本与逐次结算}
记$q_k^{(0)}$为0:00计划，$q_k^{(v)}$为6:00、12:00、18:00之后的确认量，
$v=1,2,3$。已开始时隙冻结，未开始时隙仅在发布后可调；中间电池控制无权改合同。
每次交易相对上一版记录
\begin{equation}
 \delta_{v,k}^+=\max(q_k^{(v)}-q_k^{(v-1)},0),\quad
 \delta_{v,k}^-=\max(q_k^{(v-1)}-q_k^{(v)},0).
\end{equation}
设$p_k$为执行区间固定电价，$p_v^{\rm trade}$为发布时刻所在区间价格，$U_k$为
实际紧急购电，则实际费用为
\begin{align}
 C_{\rm plan}&=\sum_kp_kq_k^{(0)},\\
 C_{\rm adj}&=\sum_{v=1}^3\sum_kp_v^{\rm trade}(1.5\delta_{v,k}^++0.5\delta_{v,k}^-),\\
 C_{\rm emg}&=\sum_k5p_kU_k,\quad
 C_{\rm total}=C_{\rm plan}+C_{\rm adj}+C_{\rm emg}.
\end{align}
采用逐版结算的明确解释：原计划费只收一次，合同未利用仍付费，减购不另退款，
不引入售电收入。先减少后恢复也要保留两笔调整，不能以最终净差抵销交易。
发布价格按所在执行区间确定，例如6:00使用$[6{:}00,6{:}10)$的价格，不错移一格。

\subsection{滚动MILP与实际执行}
每十分钟求解最长24小时的MILP，只执行第一格电池动作。目标为窗口内新增计划费、
逐次调整费、预测紧急费和次日前瞻购电代理成本减去$v_tE_{\rm end}$。
已支付费用视为常数；次日前瞻不提交合同、不入交易账本，次日0:00另建版本0。
非年度末窗口取$v_t=0.9$乘以窗口结束后144格固定电价均值。到达年度终点时，
窗口在2026年1月1日0:00截断、残值取0并保留$E_{\rm end}=6000$硬等式。
重叠窗口目标不累加为实际费用，全年费用只从执行账本重算。

窗口中的供需及储能约束为
\begin{align}
 q_k+P_k^{\rm use}+D_k+U_k&=\widehat L_k+C_k+S_k^{\rm grid},\\
 \widehat P_k&=P_k^{\rm use}+S_k^{\rm PV},\quad0\le S_k^{\rm grid}\le q_k,\\
 E_{k+1}&=E_k+0.9C_k-D_k/0.9,\quad1200\le E_k\le10800,\\
 0\le C_k&\le(5000/6)z_k,\quad0\le D_k\le(5000/6)(1-z_k),\quad z_k\in\{0,1\}.
\end{align}
本段负载与光伏已换算为kWh，充放电定义在母线侧。所有电量变量非负，且显式约束
$U_k>0\Rightarrow C_k=S_k^{\rm grid}=S_k^{\rm PV}=0$，防止用紧急电充电或同时弃电。

回放使用当格真实供需：不足时仅向下削减计划充电，仍不足才紧急购电；过剩时依次
记合同电未利用、向下削减计划放电、最后弃光。始终满足
$0\le C_k^{\rm exec}\le C_k^{\rm plan}$、$0\le D_k^{\rm exec}\le D_k^{\rm plan}$，
不增加动作、不反转方向、不在回放中重新优化。
全年诊断曾暴露实际偏差与冻结合同造成的年末不可达。经明确批准，增加只适用于
问题三的保守储备：12月31日0:00至年度终点的全部窗口状态边界均满足$E\ge6000$，
包括跨日前瞻和初态；此前下界1200、上界10800和实际年度等式不变。
该储备是额外建模假设，不是题面要求，也不保证任意偏差下的实际跨入边界或终点等式。

\subsection{指定日期与全年结果}
一月待机6000 kWh作为历史初始化，从2月1日起连续执行334天、48096格，不每日重置
SOC或施加日首尾相等。正式展示来源为明确选定的运行：
\par\noindent\texttt{q3-formal-delivery-20260913}\par
由已独立审核的全年轨迹和修正结算证据生成，不是另一次全年求解。四个指定日期的
购电展示使用最终确认量，0:00原计划另保留在\texttt{result3.xlsx}计划表中。
\IfFileExists{generated/q3_result_tables.tex}{\input{generated/q3_result_tables}}{
 \textbf{问题三同源结果文件未就绪，本文不是完整交付稿。}}

\subsection{是否引入其他时刻的预报}
相对于只用0:00预报，应把6:00、12:00、18:00的已发布版本纳入信息集，以更新未来
供需估计。但获得新预报不等于必须调整合同：只有预计减少的紧急费及改善的储能调度
足以抵偿逐次调整费时，才有经济上的调整理由。模型可以维持原合同而仅调整电池动作，
也可以在发布时刻增减未来合同，避免机械追随每一条新曲线。
发布时刻交易统计包含已计费的微小调整，说明完整版本和费用链必须保留；
该计数不是对0:00单版本策略的节费证明。本文未执行全年反事实对照，不能凭预测误差
或调整次数声称确定的节费幅度。四个官方发布时间之外没有附件预报或额外调整权，
本方案不虚构新发布；扩展必须另取得预报来源并讨论交易权限。

\subsection{检验与适用边界}
333个跨日边界的SOC完全连续，实际年度最终SOC=6000 kWh，末日145个状态边界均
不低于6000。全部区间通过功率/SOC边界、充放电互斥、供需与PV归属、禁止紧急充电
及真实输入逐格一致性检查；7717800次原始预测引用均满足因果可见性。
数值比较保持原物理容差，无事后clip。逐笔复核曾发现sidecar漏写2455笔微小已计费
调整，原汇总正确；修复只补全非零交易，原运行和审计失败证据留存，未改变费用公式。
官方模板表头与内部网格有字面冲突，导出按明确批准列序填写、保留原标签，并记录
48096格真实区间到目标cell的映射，不意味着官方已勘误。每个窗口求得最优解也不
表示实际全年策略全局最优；预测误差、受限合同和储备假设必须保留为结果解释边界。
